% Transcription — fonds Alexandre Grothendieck, shelfmark 161-3, batch 2 (pages 21-40).
% Established from the facsimile archives/batches/161-3/batch-02.pdf (a mirror of
% https://grothendieck.umontpellier.fr/161-3.pdf), per the skill
% transcribe-grothendieck. Pages 22, 27, 29 and 33 are blank — absent.
% Pass: Opus 5 (claude-opus-5), 2026-08-16 — first pass, unchecked against the pages by a human.
\documentclass[11pt,a4paper]{article}
% Le préambule commun à toutes les transcriptions du fonds Grothendieck.
%
% Il tient en une idée : **l'incertitude se compose, elle ne se résout pas.**
% Une transcription qui lisse un mot douteux a détruit la seule chose pour
% laquelle on la faisait — un lecteur doit pouvoir savoir, à chaque endroit,
% ce qui était sur le feuillet et ce qui a été ajouté. Les six macros qui
% suivent sont donc l'appareil critique, et non de la décoration.
%
% Les mêmes commandes sont reconnues par `scripts/render.mjs`, qui produit la
% vue de lecture du site. Une macro ajoutée ici doit l'être aussi là-bas,
% faute de quoi le rendu échouera bruyamment — ce qui est voulu : un rendu
% silencieusement faux est pire qu'un rendu absent.

\ProvidesPackage{grothendieck}[2026/08/08 Transcriptions du fonds Grothendieck]

\usepackage{amsmath,amssymb,amsthm}
\usepackage{mathrsfs}
\usepackage{stmaryrd}
% Les diagrammes commutatifs sont partout dans ce fonds : c'est la langue
% même de la géométrie algébrique de Grothendieck, et une transcription qui
% les rendrait en prose ne transcrirait rien.
\usepackage{tikz-cd}
% Français par défaut — c'est la langue des feuillets — mais l'anglais est
% chargé aussi : la traduction et le résumé sont des documents anglais, et la
% typographie française (espaces avant les deux-points) y serait une faute.
% Ces documents basculent par \selectlanguage{english} en tête de corps.
\usepackage[english,french]{babel}
\usepackage{fontspec}
\usepackage{geometry}
\geometry{a4paper,margin=25mm,marginparwidth=28mm}
\usepackage{marginnote}
\usepackage{xcolor}
% ulem, not soul. `soul`'s \st and \ul are built on a character-by-character
% scan that XeLaTeX's font machinery breaks: the document fails with an
% undefined control sequence rather than degrading. ulem does the same two jobs
% and is Unicode-engine safe. `normalem` stops it hijacking \emph, which the
% transcriptions use for Grothendieck's own emphasis.
\usepackage[normalem]{ulem}
\usepackage{hyperref}
\hypersetup{colorlinks=true,linkcolor=black,urlcolor=black,pdfborder={0 0 0}}

% --- Métadonnées du lot -----------------------------------------------------
% Reprises telles quelles par le script de rendu, qui les affiche en tête de
% la vue de lecture. Elles fixent ce que le fichier prétend couvrir : sans
% elles, une transcription flotte sans point d'attache dans un fonds de seize
% mille feuillets.

\newcommand{\folder}[1]{\gdef\grothFolder{#1}}
\newcommand{\batch}[1]{\gdef\grothBatch{#1}}
\newcommand{\leaves}[2]{\gdef\grothFirst{#1}\gdef\grothLast{#2}}
\newcommand{\foldertitle}[1]{\gdef\grothFoldertitle{#1}}
\gdef\grothFolder{?}\gdef\grothBatch{?}\gdef\grothFirst{?}\gdef\grothLast{?}\gdef\grothFoldertitle{}

% --- Le résumé --------------------------------------------------------------
%
% Il ouvre la lecture modernisée, sous le titre « Résumé », et il est composé
% en petit. Ce n'est pas un ornement : c'est le seul passage du document qui
% ne soit pas des mathématiques, et un lecteur doit pouvoir le reconnaître
% avant de l'avoir lu — puis le sauter s'il connaît déjà le sujet, ou s'y
% arrêter s'il ne le connaît pas. Le filet qui le suit marque la reprise du
% corps.

\newenvironment{resume}
  {\small\setlength{\parskip}{0.45em}}
  {\par\medskip\hrule\medskip}

% --- Mots-clés ---------------------------------------------------------------
%
% En anglais, en clôture du résumé de la lecture modernisée : le vocabulaire
% moderne sous lequel chercher ce que les feuillets construisent. Le manifeste
% du site (`npm run manifest`) les extrait pour étiqueter la cote entière —
% cette ligne est la source unique des tags, il n'y a pas de fichier à côté.
\newcommand{\keywords}[1]{\par\smallskip\noindent
  {\footnotesize\textsc{Keywords} --- \textit{#1}}\par}

% --- L'appareil critique ----------------------------------------------------

% Le numéro de feuillet, en marge. C'est l'ancre qui permet de régler un
% désaccord sur une formule en regardant *un* feuillet plutôt qu'une chemise
% de six cents. Le site s'en sert aussi pour tourner les pages du fac-similé
% quand on fait défiler la transcription.
\newcommand{\leaf}[1]{%
  \marginnote{\footnotesize\color{gray!70}\textsf{#1}}%
  \label{leaf:#1}\ignorespaces}

% Illisible : on ne devine pas. Un mot inventé qui se lit comme les autres est
% le pire résultat possible.
\newcommand{\ill}{\textcolor{red!60!black}{[\,\ldots\,]}}

% Lecture douteuse : le mot est donné, et signalé comme incertain.
\newcommand{\uncertain}[1]{\uline{#1}}

% Ajout de l'éditeur — une lettre manquante, un mot sous-entendu.
\newcommand{\add}[1]{\textcolor{blue!50!black}{[#1]}}

% Biffé par Grothendieck lui-même. Ce qu'il a rayé fait partie du document :
% une rature dit souvent où la pensée a changé de direction.
\newcommand{\struck}[1]{\sout{#1}}

% Note du transcripteur, sur l'état matériel du feuillet.
\newcommand{\note}[1]{\par\smallskip{\small\color{gray!60!black}\textit{[#1]}}\par\smallskip}

% Note marginale de Grothendieck, distincte du corps du texte.
\newcommand{\marginal}[1]{\par\smallskip\noindent\hspace{1em}%
  {\small\color{gray!60!black}#1}\par\smallskip}

% --- Titre ------------------------------------------------------------------

\AtBeginDocument{%
  \begin{center}
    {\large\bfseries Fonds Alexandre Grothendieck --- cote n\textsuperscript{o}~\grothFolder}\\[2pt]
    {\itshape\grothFoldertitle}\\[4pt]
    {\small feuillets \grothFirst--\grothLast\ (lot \grothBatch)}\\[2pt]
    {\footnotesize\color{gray!60!black}Transcription établie d'après le fac-similé numérisé
     par l'Université de Montpellier.\\
     Les lectures incertaines sont soulignées, les illisibles notées [\,\ldots\,],
     les ajouts entre crochets.}
  \end{center}
  \vspace{1em}}

\endinput


\folder{161-3}
\batch{2}
\pages{21}{40}
\dating{[vers 1963-1973]}
\watermark{Édition de démonstration}
\foldertitle{Topos : notes manuscrites (s.d.)}

\begin{document}

\page{21}
et soit $F \xrightarrow{f} G$ dans $\widehat{C}$, bicouvrant, prouvons que
$\widehat{u}_{0}(F) \xrightarrow{\widehat{u}_{0}(f)} \widehat{u}_{0}(G)$ est
\struck{bicouvrant} un isom. On a \ldots
\note{la page s'arrête là ; la ligne suivante est entièrement raturée. C'est la
fin de la question laissée ouverte au bas de la p.~20}

\section*{Espaces essentiels et ensembles ordonnés}

\page{23}
$X$ un ensemble

\begin{enumerate}
\item[a)] $\mathcal{T}(X)$ l'ens. des topologies sur $X$, $\underline{\omega}(X)$
l'ens. des préordres sur $X$
\[
\alpha_{X} : \mathcal{T}(X) \longrightarrow \underline{\omega}(X), \qquad
\beta_{X} : \mathcal{T}(X) \longleftarrow \underline{\omega}(X)
\]
$\alpha$ et $\beta$ fonctoriels en l'ensemble $X$ (pour images inverses de top.
resp. de relations d'ordre)

$\bullet$ à $T \in \mathcal{T}(X)$ on associe (déf.) la relation $\alpha_{X}(T)$ :
\[
y \underset{T}{\leq} x \iff x \in \overline{\{y\}}
\]

$\bullet$ à la relation $\omega$ $(\leq)$ est associé $T$ :
\[
\struck{\mathrm{Fer}}\ \mathrm{Ouv}(T) = \{\, U \subset X \mid (x \in U,\ x \leq y)
\Rightarrow y \in U \,\}
\]
$=$ réunions quelconques d'ensembles de la forme
$O_{x} = [x,\rightarrow[\ = \{ y \in X \mid y \underset{\omega}{\geq} x\}$
\[
\mathrm{Fer}(T) = \{\, Z \subset X \mid (x \in Z,\ y \leq x) \Rightarrow y \in Z \,\}
\]

\item[b)] Soit $f : X \to X'$ une application,
\struck{$X$ muni de $(T,\omega)$, $X'$ de \ill{}}

si $X, X'$ sont des espaces top. et si $f$ est \struck{continue}, $f$ est
croissante

si $X, X'$ sont préordonnés et si $f$ est croissante, $f$ est continue

\item[c)] \struck{\ill{}} D'où foncteurs

\begin{tikzcd}
(\mathrm{Esp}) \arrow[r, bend right=20, "\alpha"'] & (\text{Préord})
\arrow[l, bend right=20, "\beta"']
\end{tikzcd}
\note{sur la page les deux flèches sont superposées, $\beta$ au-dessus allant
vers $(\mathrm{Esp})$ et $\alpha$ au-dessous allant vers
$(\text{Préord})$}

\item[c)] \struck{\ill{}} Soit $T \in \mathcal{T}(X)$, montrer que tout ouvert de
$X$ est stable par générisation (resp. tout fermé de $X$ stable par
spécialisation), i.e. ouvert pour $\alpha_{X}(T)$, i.e.
\[
T \leq \beta_{X}\alpha_{X}(T)
\]
i.e. une application continue
\[
\beta\alpha(X) \longrightarrow X
\]
dont l'application sous-jacente est l'identité, d'où un homomorphisme fonctoriel
\[
\beta\alpha \xrightarrow{\ \rho\ } \mathrm{id}_{(\mathrm{Esp})}
\]
\note{le second alinéa de la page porte lui aussi la lettre c), la première ayant
été biffée puis récrite}

Soit $\omega \in \underline{\omega}(X)$, \struck{\ill{}} $T = \beta_{X}(\omega)$,
alors $x \in \overline{\{y\}}$ ssi \ill{} $y$, i.e. ssi $y \geq x$, le plus petit
ouvert $O_{x}$ contenant $x$ contient $y$, ce qui donne
\[
\alpha_{X}\beta_{X}(\omega) = \omega
\]
i.e. un isomorphisme
\[
X \xrightarrow[\ \sigma\ ]{\ \sim\ } \alpha\beta(X)
\]
tel que l'application ensembliste sous-jacente soit l'identité, d'où
\[
\mathrm{id}_{(\text{Préord})} \xrightarrow[\ \sigma\ ]{\ \sim\ } \alpha\beta
\]

\item[d)] Montrer que $\rho$ et $\sigma$ sont des morphismes d'adjonction,
\ill{} établissant pour $X \in \mathrm{Ob}(\mathrm{Esp})$,
$I \in \mathrm{Ob}(\text{Préord})$
\end{enumerate}

\page{24}
une bijection
\[
(*) \qquad \operatorname{Hom}_{\mathrm{Esp}}(\beta I, X) \simeq
\operatorname{Hom}_{\text{Préord}}(I, \alpha X)
\]

\marginal{car le foncteur $(\mathrm{Esp}) \to (\text{Préord})$ commute
aux $\varprojlim$, le foncteur $\beta$ aux $\varinjlim$. Mais de plus $\beta$ est
exact à gauche : sur $X \times Y$ la top. associée au préordre produit est le
produit des topologies. (Mais $\beta$ ne commute pas aux produits quelconques,
i.e. un produit d'espaces topologiques essentiels \ill{} essentiel.
Contre-exemple \ill{})}
\note{ce bloc est porté en diagonale dans la marge gauche, et la question qu'il
laisse ouverte est reprise à la p.~28}

En effet, une application $I \xrightarrow{f} X$ est croissante ssi elle est
continue ($X$ étant muni de la relation d'ordre $\alpha_{X}(T_{X})$, $I$ de la
topologie $\beta_{I}\omega_{I}$). En effet, si $f$ est continue, on sait qu'elle
est croissante pour les \ill{} $\alpha_{X}(T_{X})$ \ill{} ; \struck{si $f$ est
croissante} l'image inverse d'une partie de $X$ \ill{} $\beta_{X}\alpha_{X}(T_{X})$
est ouverte dans $I$, mais toute partie de $X$ (si muni de
$\beta_{X}\alpha_{X}(T_{X})$) \ill{}, donc a fortiori.
\note{le raisonnement est ici surchargé de repentirs ; seule la charpente est
lisible}

\begin{enumerate}
\item[e)] \textbf{Corollaire.} Le foncteur
\[
\beta : (\text{Préord}) \longrightarrow (\mathrm{Esp})
\]
est pleinement fidèle. Son image essentielle est formée des espaces
$(X, T_{X})$ pour lesquels $T_{X} \in \beta_{X}(\underline{\omega}(X))$ ; ce sont
\struck{celles} les espaces topologiques pour lesquels tout point $x \in X$ a un
plus petit voisinage ouvert $O_{x}$. [C'est évidemment nécessaire ; il faut
prouver que c'est suffisant, i.e. que
\[
O_{x} = \{ y \in X \mid y \underset{T}{\leq} x \}.
\]
Mais \struck{\ill{}} notons que $x \in O_{y}$ \ill{} (on cherche le plus petit
voisinage ouvert de $x$) ; $x \in \overline{\{y\}}$ signifie évidemment \ill{}
$O_{x} \ni y$, ok.]
\marginal{NB Les topologies obtenues sont aussi celles pour lesquelles toute
intersection d'ouverts est ouverte [NB ce n'est pas stable par \ill{}]}

\item[f)] Soit \struck{\ill{}} $X, T$, alors $X$ est \uncertain{sobre}
$(\overline{\{x\}} = \overline{\{y\}} \Rightarrow x = y)$ ssi le préordre
$\omega(T)$ est un ordre. \struck{\ill{}}

\item[g)] \struck{Soit $X$ muni d'un préordre, puis la topologie associée.
Désignons \ill{}} Soit $X$ muni d'un préordre, $\underline{X}$ la catégorie
associée à $X$ ordonné
$\bigl( \operatorname{Hom}(x,y) = \{\emptyset$ si $x \not\leq y$ ;
$\{e\}$ si $x \leq y\} \bigr)$, i.e.
$\operatorname{Hom}(x,y) \neq \emptyset$ ssi $O_{y} \supseteq O_{x}$. Montrer que
le foncteur $F \mapsto (x \mapsto F(O_{x}))$ de $\mathrm{Top}(X)$ dans
$\widehat{\underline{X}}$ est une équivalence de cat.
\[
\mathrm{Top}(X) \simeq \widehat{\underline{X}}
\]
On identifie les $a \subset X$ à des sommes \ill{} des $e_{\underline{X}}^{\wedge}$,
et les cribles de $X$ (i.e. les ouverts de $X$) à des sous-objets de
$e_{\underline{X}}^{\wedge}$, i.e. \ill{}
\end{enumerate}

\page{25}
\begin{enumerate}
\item[g)] $X$ muni d'une topologie $T$ déduite d'un préordre $\omega$ [dans
laquelle les ouverts \ill{} les voisinages, sont de la forme $O_{x}$, $x \in X$].
$X$ est irréductible ssi deux ouverts non vides ont une intersection non vide,
i.e. pour $x, y \in X$, $O_{x} \cap O_{y} \neq \emptyset$, i.e.
$\forall x,y \in X$, $\exists z \in X$ avec $z \leq x$, $z \leq y$, i.e. $X$ est
\emph{cofiltrant}. En conclure que les fermés irréductibles de $X$ sont les
parties $Y$ de $X$ telles que
\begin{enumerate}
\item[a)] $y \in Y$, $y' \leq y \Rightarrow y' \in Y$ \quad ($Y$ fermé)
\item[b)] $\forall x,y \in Y$, $\exists z \in Y$, tel que $z \leq x$, $z \leq y$.
\end{enumerate}
\note{cette lettre g) est la seconde de la suite : la p.~24 en portait déjà une,
celle qui introduit $\underline{X}$ et à laquelle renvoie le h) ci-dessous}

\item[h)] $X$ comme dans g), $\underline{X}$ la catégorie associée à la relation
de préordre, car $\operatorname{Hom}_{\underline{X}}(x,y) \leq 1$,
$\operatorname{Hom}_{\underline{X}}(x,y) \neq \emptyset \Rightarrow x \leq y
\iff y \in \overline{\{x\}}$.

Considérons le foncteur $\mathrm{Top}(X) \to \widehat{\underline{X}}$,
$F \mapsto (x \mapsto F(O_{x}))$. Prouver que ce foncteur est une équivalence
\[
\mathrm{Top}(X) \simeq \widehat{\underline{X}}
\]
Dorénavant on identifie [\struck{si} $X$ précisant i.e. \emph{ordonné}] $X$ à un
ensemble de sous-objets de l'objet final de
$\widehat{\underline{X}} = \mathrm{Top}(X)$ [\struck{cribles} \ill{} des topos],
et de \ill{} par la \ill{} de $X$, \ill{} les sous-objets de
$e_{\underline{X}}$ \ill{} aux ouverts des topos.

Interpréter
\[
\mathrm{Points}(\widehat{\underline{X}}) \simeq \mathrm{Pro}(\underline{X})
\qquad (\ill{}\ \mathrm{ind}),
\]
en identifiant les pro-objets de $\underline{X}$ aux sous-ensembles ordonnés $Y$
de $X$ satisfaisant les conditions a) b) de g), \struck{et} la structure de
catégorie provenant de la relation d'ordre entre parties de $X$ \emph{opposée} à
la relation d'inclusion. Soit \struck{$\mathrm{Pro}(X)$} $\widetilde{X}$
l'ensemble ordonné ainsi construit. L'application
\[
x \longmapsto \overline{\{x\}} = \{ y \in X \mid y \geq x \}
\]
est une application $X \to \widetilde{X}$, qui s'identifie à l'application
canonique de $X$ dans l'espace sobre associé, ou à
$\underline{X} \to \mathrm{Pro}(\underline{X})$ (foncteur canonique).

\item[i)] Conditions équivalentes
\begin{enumerate}
\item[$\alpha$)] $X \to \widetilde{X}$ surjective
\item[$\beta$)] Toute partie fermée cofiltrante de $X$ a un plus petit élément.
\end{enumerate}
\end{enumerate}

\page{26}
\begin{enumerate}
\item[$\gamma$)] \struck{$T_{0}$} $X$ est \struck{(pour sa relation d'ordre)}
artinien, i.e. toute \struck{famille} suite décroissante d'éléments est
stationnaire.
\item[$\delta$)] $\widetilde{X}$ est essentiel (i.e. chaque
\uncertain{point} de $X$ est essentiel)
\end{enumerate}
\marginal{en fait, $\operatorname{Im}(X \to \widetilde{X})$ est l'ens. des
\ill{} ess. de $\widetilde{X}$ !}

\begin{enumerate}
\item[j)] Soit $E$ un topos. Conditions équivalentes
\begin{enumerate}
\item[$\alpha$)] $E$ \struck{n'est} est équivalent à un $\mathrm{Top}(X)$ et à un
$\widehat{C}$
\item[$\beta$)] $E$ équiv. à $\mathrm{Top}(X)$, $X$ espace essentiel
\item[$\gamma$)] $E$ équiv. à $\widehat{I}$, $I$ ens. ordonné
\item[$\delta$)] $E$ engendré par les sous-objets de l'objet final, et a
suffisamment de points essentiels
\end{enumerate}
[On dit alors que $E$ est un espace top. \emph{essentiel}.]
\end{enumerate}

Prouver que si $C$ est une catégorie, $\widehat{C}$ satisfait aux conditions
précédentes ssi $C$ est préordonnée.

Prouver que si $X$ est un espace top., $\mathrm{Top}(X)$ satisfait aux conditions
précédentes ssi l'ens. des $x \in X$ qui sont essentiels est très dense.

\struck{\ill{}} On a des équivalences de $2$-catégories
\[
(\text{topos spatiaux essentiels}) \approx (\text{Préord})
\ \struck{\text{ensembles ordonn\'es}}\ \approx (\mathrm{Esp.\ ess.})
\]
[NB. Si $X$ est essentiel on le \ill{} à partir de $E = \mathrm{Top}(X)$, comme
\ill{} $\mathrm{Points}(E)$.]

\page{28}
\struck{$\gamma$} commute aux $\inf$ finis ?
\[
\gamma(X \times Y) = \alpha(X) \wedge \beta(Y) = 1_{\sigma}
\]
\begin{align*}
\gamma(W \cap W') &= \sup_{U \times V \subset W \cap W'} \alpha(U) \wedge \beta(V)
= \sup_{\substack{U \times V \subset W \\ U' \times V' \subset W'}}
\alpha(U \cap U') \wedge \beta(V \cap V') \\
&= \sup_{U \times V \subset W} \alpha(U)\beta(V) \wedge
\sup_{U' \times V' \subset W'} \alpha(U')\beta(V') \\
&= \gamma(W) \wedge \gamma(W')
\end{align*}

$X \times Y$

\begin{tikzcd}[column sep=large]
O_{X} \arrow[r, "{U \mapsto (U,V)}"] & O_{X} \times O_{Y}
\arrow[r, "{(U,V) \mapsto U \times V}"] & \underline{O}_{X \times Y} \\
O_{Y} \arrow[ur, "{V \mapsto (X,V)}"'] & &
\end{tikzcd}

\marginal{$\sigma$ \ill{} et $\sup$ quelconques, et $\inf$ finis distributifs
par rapport aux $\sup$ quelc.}
\marginal{$\operatorname{Hom}'$ \ill{} i.e. \ill{} commutant aux $\sup$
quelconques et $\inf$ finis}

\[
u \longmapsto \alpha, \beta \qquad
\operatorname{Hom}'(\underline{O}_{X \times Y}, \sigma) \longrightarrow
\operatorname{Hom}'(O_{X}, \sigma) \times \operatorname{Hom}'(O_{Y}, \sigma)
\]
injectif (car $u \in \operatorname{Hom}'(\underline{O}_{X \times Y}, \sigma)$ est
connu dès que l'on connaît $u(U \times V) = \alpha(U) \wedge \beta(V)$, $\gamma$
étant \ill{}, puisque tout $W \in \underline{O}_{X \times Y}$ est sup de tels
$U \times V$).

Est-ce bijectif ? soit $\alpha : O_{X} \to \sigma$, $\beta : O_{Y} \to \sigma$,
posons
\[
\gamma(W) = \sup_{\substack{U \in O_{X},\ V \in O_{Y} \\ U \times V \subset W}}
\alpha(U) \wedge \beta(V)
\]
(c'est une application croissante
$\underline{O}_{X \times Y} \to \sigma$). \quad OBS $X \neq \emptyset$,
$Y \neq \emptyset$

$\alpha(X' \times Y) = ?$ Notons que si $Y \neq \emptyset$, alors
$U \times V \subset X' \times Y$ ssi $U \subset X'$, donc le sup est
\struck{égal} $\sup \alpha(U) \wedge \beta(V) = \alpha(X')$, où
$X' \in \mathrm{Ouv}(X)$.

\note{la moitié supérieure droite de la feuille porte un calcul sans rapport avec
ce qui précède — quaternions et formules de duplication trigonométriques,
$z = e^{i\theta}$, $u = a + bi + cj$, $a^{2}+b^{2}+c^{2}=1$ — qui n'est pas
transcrit ici}

\section*{Lemme sur les modules localement libres et les algèbres d'Azumaya}

\page{30}
\textbf{Lemme.} $E$ un topos \struck{annelé} localement annelé (dans les deux
cas)

\begin{itemize}
\item $C$ catégorie des couples $(\mathcal{E}, \mathcal{G})$ de modules loc.
libres de t.f. partout $\neq 0$
\item $C'$ \rule{2em}{0.4pt} quadruples $(A, \mathcal{E}, \mathcal{F}, \varphi)$,
avec
$A$ Algèbre sur $E$ loc. isom. à une Alg. de matrices ; $\mathcal{E}, \mathcal{F}$
Mod. loc. lib. de t.f. ; $\varphi$ un isom.
$A \otimes \underline{\operatorname{End}}(\mathcal{E}) \simeq
\underline{\operatorname{End}}(\mathcal{F})$
\end{itemize}
\marginal{$A$, $\mathcal{E}$, $\mathcal{F}$ partout $\neq 0$}

$F : C \to C'$ le foncteur
\[
F : (\mathcal{G}, \mathcal{E}) \longmapsto
\underline{\operatorname{End}}(\mathcal{G}),\ \mathcal{E},\
\mathcal{G} \otimes \mathcal{E},\ \varphi_{\mathcal{G},\mathcal{E}}
\]
où $\varphi_{\mathcal{E},\mathcal{G}}$ est l'isom. canonique
$\underline{\operatorname{End}}(\mathcal{G}) \otimes
\underline{\operatorname{End}}(\mathcal{E}) \simeq
\underline{\operatorname{End}}(\mathcal{G} \otimes \mathcal{E})$.

Alors $F$ est une équivalence.

\emph{Dém.} $F$ provient d'un hom. de champs $\underline{C} \to \underline{C}'$,
et il suffit de prouver que celui-ci est une équivalence. Or, si
$\mathcal{G}, \mathcal{E}$ (resp. $A, \mathcal{E}, \mathcal{F}$ tous de rang
$\geq 0$) sont des entiers $> 0$, soit $\underline{C}_{g,e}$ le sous-champ de
$\underline{C}$ formé des $\mathcal{E}, \mathcal{G}$ tels que
$\operatorname{rang}\mathcal{E} = e$, $\operatorname{rang}\mathcal{G} = g$
(resp. $\underline{C}'_{a,e,f}$ le sous-champ de $\underline{C}'$ formé des
quadruples pour lesquels $A$ de rang $a^{2}$, $\mathcal{E}$ de rang $e$,
$\mathcal{F}$ de rang $f$). On a
\[
\underline{C} \simeq \coprod_{g,e} \underline{C}_{g,e} \qquad
\underline{C}' \simeq \coprod_{a,e,f} \underline{C}'_{a,e,f}
\]
(somme de champs) et $F$ induit
\[
F_{g,e} : \underline{C}_{g,e} \longrightarrow \underline{C}'_{g^{2},e,\ill{}}
\]
on est donc ramené à prouver que les $F_{g,e}$ sont des équivalences.

Notons \struck{$\underline{C}_{g,e}$ est une gerbe (liée par
$Gl(e) \times Gl(g) = \operatorname{Aut}(\mathcal{G})$)} \struck{\ill{}}
\note{les deux lignes qui suivent sont entièrement biffées ; c'est le mot
« gerbe » qui est abandonné}

\struck{En effet, $\underline{C}_{g,e}$ loc. \ill{} a un \ill{}} En effet, c'est
essentiellement un champ en groupoïdes, donc reste à prouver que deux objets sont
loc. isom. On peut supposer $A = A'$, $\mathcal{E} = \mathcal{E}'$,
$\mathcal{F} = \mathcal{F}'$. On considère deux isom.
\[
\varphi, \varphi' : A \otimes \underline{\operatorname{End}}(\mathcal{E})
\rightrightarrows \underline{\operatorname{End}}(\mathcal{F})
\]
Ils diffèrent par un automorphisme \struck{intérieur} de
$\underline{\operatorname{End}}(\mathcal{F})$, d'où un automorphisme de
$\mathcal{F}$ qui est (localement) intérieur, donc on \ill{}. \struck{De plus}

Pour terminer la démonstration, il reste à prouver que $F_{g,e}$ induit
$\operatorname{Aut}(\xi) \xrightarrow{\alpha} \operatorname{Aut}(\varphi(\xi))$
pour \ill{} un isom. Or $\xi = (\mathcal{G}, \mathcal{E})$,
$\xi' = \underline{\operatorname{End}}(\mathcal{G}), \mathcal{E},
\mathcal{G} \otimes \mathcal{E}, \varphi_{\mathcal{G},\mathcal{E}} :
\underline{\operatorname{End}}(\mathcal{G}) \otimes
\underline{\operatorname{End}}(\mathcal{E}) \simeq
\underline{\operatorname{End}}(\mathcal{G} \otimes \mathcal{E})$, et si $u, v$ est
un automorphisme de $\xi$, $F(u,v)$ est l'automorphisme
$\underline{\operatorname{End}}(u), v, u \otimes v$, donc si $F(u,v)$ est connu,
$v$ est connu et $u \otimes v$ est connu, d'où facilement (comme $\mathcal{E}$
partout $\neq 0$) $u$ connu, donc $\alpha$ injectif. Soit maintenant
$\tau = (\lambda, v, w)$ un automorphisme

\page{31}
\note{page au crayon, restée à l'état d'esquisse ; on n'en transcrit que ce qui
est formé}

\begin{tikzcd}
C^{1} \arrow[r] & C^{2} \arrow[r] & C^{3} \arrow[r] & C^{4}
\end{tikzcd}

$C = \underline{\operatorname{Hom}}(K, K')$

\[
\mathrm{TORS}^{1}_{U}(X, \Pi) \simeq H^{1}(X, \Pi)
\]

\begin{tikzcd}[column sep=small, row sep=small, nodes={font=\scriptsize}]
{[C^{0}]} \arrow[r] & C^{0} \times Z^{1} \arrow[d, bend right=20]
\arrow[d, bend left=20] & \\
& Z^{1} \arrow[r] & C^{1} \times Z^{2} \arrow[d, bend right=20]
\arrow[d, bend left=20] \\
& \operatorname{Aut}(0) \arrow[u] & 0 \to Z^{2}
\end{tikzcd}

\[
T^{n} = \mathrm{TORS}^{n}_{U}(X ; \Pi) \rightleftarrows
\uncertain{H\!I\!I^{n}}(X, \Pi) \implies \pi_{0}(T^{n}) \simeq H^{n}(X, \Pi)
\]
\[
\mathrm{TORS}^{1}_{U}(X, \Pi) \xrightarrow{\ \sim\ } \uncertain{H\!I^{1}}(X, \Pi)
\]

\begin{tikzcd}[column sep=small, row sep=small]
E_{\bullet} \arrow[d, bend right=20] \arrow[d, bend left=20] &
\Pi \times E \arrow[d, bend right=20] \arrow[d, bend left=20] & \Pi \arrow[d] \\
E_{1} \arrow[d, bend right=20] \arrow[d, bend left=20] & E \arrow[d] & 1 \\
E_{0} \arrow[d] & X & \\
X & &
\end{tikzcd}

\page{32}
de $\xi'$, donc $\lambda$ un automorphisme de
$\underline{\operatorname{End}}(\mathcal{G})$, $v$ de $\mathcal{E}$, $w$ de
$\mathcal{G} \otimes \mathcal{E}$, avec commutativité de

\begin{tikzcd}[column sep=large]
\underline{\operatorname{End}}(\mathcal{G}) \otimes
\underline{\operatorname{End}}(\mathcal{E})
\arrow[r, "\varphi_{\mathcal{G},\mathcal{E}}"]
\arrow[d, "\lambda \otimes \underline{\operatorname{End}}(v)"'] &
\underline{\operatorname{End}}(\mathcal{G} \otimes \mathcal{E})
\arrow[d, "\underline{\operatorname{End}}(w)"] \\
\underline{\operatorname{End}}(\mathcal{G}) \otimes
\underline{\operatorname{End}}(\mathcal{E})
\arrow[r, "\varphi_{\mathcal{G},\mathcal{E}}"] &
\underline{\operatorname{End}}(\mathcal{G} \otimes \mathcal{E})
\end{tikzcd}

On sait qu'il suffit de voir que loc. il provient d'un (unique) automorphisme de
$\xi$, qui sera de la forme $(u,v)$, $u$ un automorphisme de $\mathcal{G}$, donc
on doit prouver
\[
\lambda \simeq \underline{\operatorname{End}}(u), \qquad w = u \otimes v
\]
Or on sait que loc. $\lambda$ est \struck{\ill{}} intérieur, i.e. qu'il existe
loc. $u$ tel que $\lambda = \underline{\operatorname{End}}(u)$. Reste \ill{} en
remplaçant \ill{} par $F(u,v)^{-1}$, on peut supposer que $\lambda = \mathrm{id}$.
Donc on doit a
\[
\varphi \circ (\mathrm{id} \otimes \underline{\operatorname{End}}(v)) \simeq
\underline{\operatorname{End}}(w) \circ \varphi,
\]
et on doit prouver qu'il existe [loc.] une section $u$ de \ill{} telle que
$w = \ill{}\,(\mathrm{id} \otimes v)$, i.e.
$w = \mathrm{id}_{\mathcal{G}} \otimes v$ \ill{}

\struck{prenant la même image
$1_{g} \otimes \underline{\operatorname{End}}(\mathcal{E}) \to
\underline{\operatorname{End}}(\mathcal{G}) \otimes
\underline{\operatorname{End}}(\mathcal{E})$}, la commutativité
\struck{partielle} donne
\[
\underline{\operatorname{End}}(w)(\mathrm{id}_{g} \otimes X) =
\mathrm{id}_{g} \otimes \underline{\operatorname{End}}(v)(X) \qquad
\text{pour } X \in \Gamma\,\underline{\operatorname{End}}(\mathcal{E})
\]
\[
w(\mathrm{id}_{g} \otimes X)w^{-1} = \mathrm{id}_{g} \otimes vXv^{-1}
\]
\ill{} s'écrit que pour tout end. $X$ de $\mathcal{G}$, $Y$ de $\mathcal{E}$,
\ill{}
\note{la fin de la page est trop surchargée pour être suivie ; l'argument
s'arrête ici}

\section*{Notes de cours : topos et espaces sobres}

\page{34}
\emph{Recall}: direct and inverse images, fiber functor as inverse images, direct
images of \struck{sheaves with} algebraic structures (versus inverse images)
(defined by finite $\varprojlim$ and by $\varinjlim$) / case of an étale
morphism: its base change \ill{} side $\mathrm{Top}(X)$.

Examples: field, integral domain. Example: product of epimorphisms not an
epimorphism.

\marginal{sheaf of \ill{} ordered \ill{}. Examples. Constant sheaves.
\struck{Sheaves} $\underline{\operatorname{Hom}}(F,G)$ and $\beta(F)$. Their
intrinsic character. $\beta(F) \simeq \underline{\operatorname{Hom}}(F, \beta(\ill{}))$.
$\mathrm{Top}(X) \simeq$ \ill{} functors \ill{} which commute with $\varprojlim$}

\struck{The abstract} category $T(X)$ embodies all topological information
workable for $X$:
\begin{enumerate}
\item[a)] the notion of \struck{sheaves} with alg. structures --- $X$ can be
reconstructed from $T(X)$.
\item[b)] The notion of \emph{stack} over $X$ depends on $T(X)$ alone (up to
equivalence)
\item[c)] Connexity of $X$ ; \struck{\ill{}} cohomology and homotopy
\item[d)] Relation \struck{invariants} between $\pi_{1}$ and covering spaces of
$X$ \quad (cf. SGA 4 V and Artin \& Mazur)
\item[e)] Non commutative \struck{homological} algebra: $H^{1}(\ ,\ )$
\item[f)] Other examples: local connectedness ; total disconnectedness,
\struck{quasi}-compactness \ldots
\item[g)] Operations on spaces (sums, amalgamated sums) and properties of the
continuous maps reflected by operations on topoi.
\end{enumerate}

The « reasonable » spaces $X$ can be recovered from $\mathrm{Top}(X)$, and so can
the continuous maps.

$\mathrm{Fib}(E)$ (isomorphism classes of \emph{generalized} fiber functors) on a
topos space. It corresponds to the set $\widetilde{X}$ of all closed irred.
subsets of $X$, with the obvious topology, and

NB $\mathrm{Fib}(E)$ is rigid if $E = \mathrm{Top}(X)$,
\struck{$X \to \mathrm{Fib}(E)$}

$X \to \mathrm{Fib}(E)$ can be identified with $x \mapsto \overline{\{x\}}$. The
topology of $X$ is the inverse image topology.

map injective $\iff \overline{\{x\}} = \overline{\{y\}} \Rightarrow x = y$ (i.e.
separation axiom that $x \neq y$ implies that either $x$ or $y$ has a
neighbourhood not containing the other)

map surjective $\iff$ every closed irred. subset of $X$ has a generic point

\struck{so-called} map bijective $\iff x \mapsto \overline{\{x\}}$ is a bijective
map onto $\widetilde{X}$ : \quad $X$ \emph{sober}

[$X \to \widetilde{X}$ can be viewed as being the universal map of $X$ into a
sober space, $X \mapsto \widetilde{X}$ the left adjoint of the inclusion
$(\mathrm{Sp\ sob}) \longrightarrow (\mathrm{Sp})$
\[
\operatorname{Hom}_{\mathrm{Sp\ sob}}(X, i(Z)) \simeq
\operatorname{Hom}_{\mathrm{Sp}}(\widetilde{X}, Z)
\]
Then $X \to \widetilde{X}$ commutes with arbitrary $\varprojlim$ \ldots]
\note{les deux indices de catégorie sont portés dans cet ordre sur la page, alors
que l'adjonction énoncée juste avant les échangerait ; on n'y touche pas}

\marginal{\ill{} $\mathrm{Top}(X) \simeq \mathrm{Top}(\widetilde{X})$,
$\mathrm{Fib}(\mathrm{Top}(X))$, \ill{} it is a \ill{} bijective, it is done
\ill{}}

\page{35}
\note{la page porte, de la main de Grothendieck, le numéro 1 ; les pages 37 et 39
portent 3 et \ill{} et poursuivent la même suite}

$T$ espèce de structure « définie par $\varprojlim$ finies » (resp. quelc.) sur
famille d'ens. de base $(E_{i})_{i \in I}$ \quad ($I$ ens. donné)

Pour toute catégorie $C$ avec $\varprojlim$ finies (resp. quelconques) on a donc
\[
\text{« famille d'objets de } C \text{ sous-jacente »} \qquad
T(C) \xrightarrow{\ \varphi^{T}_{C}\ } C^{I}
\]
foncteur allant de la catégorie des « structures d'espèce $T$ dans $C$ » vers
$C^{I}$. Pour tout foncteur
\[
u : C \longrightarrow C'
\]
commutant aux $\varprojlim$ finies (resp. quelc.) on a de plus
$T(C) \xrightarrow{T(u)} T(C')$, rendant commutatif le diagramme

\begin{tikzcd}
T(C) \arrow[r, "T(u)"] \arrow[d, "\varphi^{T}_{C}"'] & T(C')
\arrow[d, "\varphi^{T}_{C'}"] \\
C^{I} \arrow[r, "u^{I}"'] & C'^{I}
\end{tikzcd}

On a $T(\mathrm{id}_{C}) = \mathrm{id}_{T(C)}$, $T(vu) = T(v)T(u)$, donc $T$ est
un foncteur de $L_{0} = (\mathrm{Cat}\ \varprojlim \text{ finies})$ (resp.
$L = (\mathrm{Cat}\ \varprojlim)$) dans $\mathrm{Cat}$, d'où \struck{on} une
catégorie \ill{}-fibrée $\mathcal{T}$ \ill{}, muni d'un \struck{foncteur} \ill{}
de $T$ dans $\mathcal{T}_{I} : C \mapsto C^{I}$,
\[
\varphi^{T} : \mathcal{T} \longrightarrow \mathcal{T}_{I}
\]
Soit $T'$ une autre « espèce de structure », alors on définit
\[
\underline{\operatorname{Hom}}_{L}(T, T') =
\underline{\operatorname{Hom}}_{L_{0}}(\mathcal{T}_{L_{0}}, \mathcal{T}'_{L_{0}})
\]
\[
\bigl[\ \operatorname{resp.}\ \underline{\operatorname{Hom}}_{L}(T, T') =
\underline{\operatorname{Hom}}_{L_{0}}(\mathcal{T}_{L_{0}},
\mathcal{T}'_{L_{0}}) \ \bigr]
\]
et si $I = I'$
\[
\underline{\operatorname{Hom}}_{L}(I ; T, T') = \text{catégorie des hom. de }
\mathcal{T}_{L_{0}} \text{ dans } \mathcal{T}'_{L_{0}}
\]
munis d'une donnée de compatibilité $\alpha$ entre $\varphi^{T}$ et
$\varphi^{T'}$ à valeurs dans $\mathcal{T}_{I}$.

\note{la variante « resp. » annoncée entre crochets au bas de la page est laissée
en blanc}

\page{36}
in other words, $X$ and $\widetilde{X}$ have the same sheaf theory. For all
(most!) practical purposes, $X$ can be replaced with $\widetilde{X}$.

We have a map
\[
\operatorname{Hom}_{\mathrm{Sp}}(X, Y) \longrightarrow
\operatorname{Hom}\mathrm{top}(\mathrm{Top}(X), \mathrm{Top}(Y))
\qquad (\text{isom classes})
\]
it is bijective if $Y$ is sober. In fact, the category
$\underline{\operatorname{Hom}}\mathrm{top}(\mathrm{Top}(X), \mathrm{Top}(Y))$ is
\struck{discrete} rigid, associated to an order relation on
$\operatorname{Hom}\mathrm{top}(\mathrm{Top}(X), \mathrm{Top}(Y))$
(specialisation).

\medskip

\textbf{The notion of a topos.} \quad (gener.) « topology » $=$ study of topoi.

(Why it is a natural and useful generalisation of the notion of topological
space?
\begin{enumerate}
\item[a)] It is the natural setting for most « algebraic » (and non algebraic)
structures.
\item[b)] As such, it allows to construct classifying topoi for such structures
(possibly all): algebra $=$ topology!
\item[c)] It is stable under a variety of operations that in many cases do not
make sense in the category of topological spaces:
$*$ generalized notion of passing to quotient (by wild equivalence relations,
which would give pathological quotients within top. spaces)
$*$ clutching together of pieces à la Artin
\item[d)] They arise in various questions of geometry \struck{as modular} topoi
\item[e)] They provide the most convenient set-up for all topics involving
systematically passing to quotients.
\item[f)] They give the ultimate outcome for [the intuitive notion of]
localization
\item[g)] They are technical tools for various cohomology theories in algebraic
geometry.
\end{enumerate}

\textbf{Examples of topoi} \quad $\mathrm{Top}(X)$, $\mathrm{Top}(X, G)$,
$\widehat{C}$ ; for instance $B_{G}$ ($G$ discrete) ;
$\mathrm{Top}(X)$, $B_{G}$, $H^{*}(G, -)$

\page{37}
\note{page numérotée 3 de la main de Grothendieck}

Si $T$ est donné avec $I$, on a une application canonique
\[
I \longrightarrow \mathrm{Ob}\,\Gamma^{T}_{\varphi} \qquad
(\text{resp. } I \to \mathrm{Ob}\,\Gamma^{T}_{\varphi_{0}})
\]
On a, si $T$ est de $\varprojlim$, un foncteur
\[
\Gamma^{T} \longrightarrow \Gamma_{0}^{T}
\]
(En particulier \ill{} pour $T' = \tau$) et on a commutativité dans

\begin{tikzcd}
\mathrm{Ob}\,\Gamma^{T} \arrow[rr] & & \mathrm{Ob}\,\Gamma_{0}^{T} \\
& I \arrow[ul] \arrow[ur] &
\end{tikzcd}

\struck{Supposons}

\textbf{Ex.} Supposons qu'il existe un
$\beta^{T} \in (\mathrm{Cat}\ \varprojlim)$ (resp.
$\beta_{0}^{T} \in (\mathrm{Cat}\ \varprojlim)$) telle que $T$ (resp. $T_{0}$)
soit \struck{défini} $2$-représentable par cet objet :
\[
T(C) \simeq \underline{\operatorname{Hom}}_{\varprojlim}(\beta^{T}, C) \qquad
C \in \mathrm{Ob}(\mathrm{Cat}\ \varprojlim)
\]
\[
\bigl(\ T(C) \simeq \underline{\operatorname{Hom}}_{\varprojlim}(\beta_{0}, C)
\qquad C \in \mathrm{Ob}(\mathrm{Cat}\ \varprojlim)\ \bigr)
\]
Alors
\begin{enumerate}
\item[a)] $\beta$ ($\beta_{0}$) est défini à équivalence près,
\struck{défini} unique à isom. unique près
\item[b)] Pour toute catégorie \uncertain{$2$-cofibrée} $\mathcal{F}$ sur
$(\mathrm{Cat}\ \varprojlim)$ (resp. $(\mathrm{Cat}\ \varprojlim)$) on a une
équiv. can.
\[
\underline{\operatorname{Hom}}(T, \mathcal{F}) \simeq \mathcal{F}(\beta)
\]
\[
\bigl(\ \underline{\operatorname{Hom}}_{0}(T, \mathcal{F}) \simeq
\mathcal{F}(\beta_{0})\ \bigr)
\]
\end{enumerate}
en particulier, si $\mathcal{F}$ est défini par $T'$ on a
\[
\underline{\operatorname{Hom}}(T, T') \simeq T'(\beta)
\]
\[
\bigl(\ \underline{\operatorname{Hom}}_{0}(T, T') \simeq T'(\beta_{0})\ \bigr)
\]
et en particulier, faisant $T' = i^{\text{struct}}$, on a $T'(\beta) = \beta$
($T'(\beta_{0}) = \beta_{0}$), on trouve
\[
\beta = \underline{\operatorname{Hom}}(T, i) = \Gamma^{T}
\]
\[
\bigl(\ \beta_{0} = \underline{\operatorname{Hom}}_{0}(T, i_{0}) =
\Gamma_{0}^{T}\ \bigr)
\]

\page{38}
\textbf{Complements on relations between $X$ and $\mathrm{Top}(X)$}

\begin{enumerate}
\item[a)] Specialisations \struck{and} morphisms between fiber functors

If $x \in \overline{\{y\}}$, then we have a functorial homomorphism
$F_{x} \to F_{y}$ ($F \in \mathrm{Ob}\,\mathrm{Top}(X)$) because the
neighbourhoods of $y$ contain $x$. Thus \struck{\ill{}} we get
\[
\mathrm{Cat}\bigl(X \text{ ordered by } x \leq y \text{ iff }
x \in \overline{\{y\}}\bigr) \longrightarrow \mathrm{Fib}(\mathrm{Top}(X))
\]

\textbf{Proposition.} $\operatorname{Hom}(\xi_{x}, \xi_{y}) =
\{\emptyset$ if $x \not\leq y$ ; reduced to the specialisation hom. otherwise$\}$,
i.e. the previous functor is \emph{fully faithful}.

\textbf{Corollary.} If $X$ sober (\struck{and only} more generally if
$X \to \widetilde{X}$ surjective, \ldots) then (and only then \ldots) the
previous functor is an equivalence. [Assume $X$ \struck{\ill{}} for simplicity]

\item[b)] $\mathrm{Cat}\text{-}\underline{\operatorname{Hom}}(X, Y)$ ordered by
$f \leq g$ iff $\forall x \in X$, $f(x) \in \overline{\{g(x)\}}$
\[
\downarrow
\]
\[
\underline{\operatorname{Hom}}(\mathrm{Top}(Y), \mathrm{Top}(X))
\]
\emph{fully faithful} \struck{\ill{}}. If $Y$ is ordered it is essentially
\ill{} consists of the $f^{*}$'s that are left exact and commute with
$\varprojlim$, and \struck{\ill{}}
\end{enumerate}

\marginal{NB in the case \ill{} of the \ill{} $2$-category \ill{} we get lots of
fiber functors \ill{} we obtain « certain topoi » \ill{}}

\medskip

\textbf{Topos.}
\begin{enumerate}
\item[1)] \textbf{Definition}
\begin{enumerate}
\item[a)] Through certain properties $+$ generators (for instance a) b) c) d))
\item[$*$] Many sets of equivalent \struck{axioms} properties possible. What is
it \struck{\ill{}} that governs us will \ill{}, whatever those properties \ldots\
comes up?

The original definition was in terms of categories of \struck{filtered}
(generalized sheaves) or categories \ill{} satisfying \ill{} the categories
\struck{\ill{}} $O_{X}$. We cannot express it now, but we can express
\item[b)] $E \xrightarrow[f_{*}]{\text{fully faithful}}$ some $\widehat{C}$,
$C$ « small », with a left adjoint $f^{*}$ which is \emph{left exact}.
\end{enumerate}
\end{enumerate}

\page{39}
\note{feuillet relié à l'envers dans le dossier ; il poursuit la suite numérotée
des pages 35 et 37}

Considérons l'espèce de structure « anneau (commutatif) », soit $T$. On sait
qu'elle est définie en termes de $\varprojlim$ finies, soit
$\lambda_{0} = (\text{ensemble fini}, \ill{})$ le type de structure
correspondant. On se propose donc d'étudier, dans plus tard, tout $\lambda$ type
tel que $\lambda_{0} \to \lambda$ (les $\lambda$-catégories \ill{} des produits
finis, les $\lambda$-foncteurs y commutent) : il définit un $\lambda$-type
$T_{\lambda}$. On se propose, pour divers $\lambda$, d'étudier l'existence
\emph{et l'unicité} d'un sous-type plein $T^{*}_{\lambda}$
($T^{*}_{\lambda}(C) \to T_{\lambda}(C)$ pl. fidèle pour tout $C$) de telle façon
que pour $C = (\mathrm{Ens})$, $T^{*}_{\lambda}(C)$ soit la sous-catégorie pleine
de $T_{\lambda}(C) = $ cat. des anneaux (comm. unifères) qui sont des
\emph{corps} (commutatifs).

[On sait que l'unicité est acquise si $\lambda(\lambda_{1}, \lambda_{2} =
\emptyset)$]
\begin{enumerate}
\item[a)] si $\lambda = \lambda_{0}$, il n'existe pas une telle
$T_{\lambda}$. En effet, si $\lambda = (\lambda_{1}, \lambda_{2} = \emptyset)$,
$\lambda$ formé de $\varinjlim$ \ill{} finis, alors $T_{\lambda}$ est
$\lambda$-représentable par la $\lambda$-catégorie $R^{\lambda}_{T}$ qui est la
sous-catégorie pleine $S^{T}_{\lambda}$ de $\mathrm{Ann}$ engendrée (par les
$\varinjlim$ de type $\lambda$) par $\mathbb{Z}[T]$ (l'objet libre engendré par
$1$ élément). Les sous-\struck{catégories} types pleins de $T_{\lambda}$
\struck{sont les} correspondent aux catégories de fractions de
$S^{T}_{\lambda}$ par des ens. de flèches $M$ stables par composition, et par
$\varinjlim$ de type $\lambda$. (On aura cat. \struck{quotient} de fractions
analogue pour $S^{T}_{\lambda}$ pour des ens. $M$ de flèches stables par
$\varinjlim$ de type $\lambda$.) Ici, les flèches en question sont \struck{\ill{}}
telles \struck{\ill{}} des flèches $u : A \to B$ dans
$S^{T}_{\lambda}$ ($\subset \mathrm{Ann}$) telles que pour tout corps
$k \in \mathrm{Ob}(\mathrm{Ann})$, la flèche
$\operatorname{Hom}(B, k) \to \operatorname{Hom}(A, k)$ soit
\emph{bijective} [ou encore que les morphismes de schémas affines
$u' : X \to \struck{Y}$ définis par $u$ soient tels que $X(k) \simeq Y(k)$ pour
tout corps $k$. On sait que cela signifie que $u$ est \struck{bijective} et
définit des ext. résiduelles triviales]. \struck{\ill{}} Cet ens. de flèches est
évidemment stable par composition et par $\varinjlim$ (resp. $\varprojlim$) de
type $\lambda$. (Serait vrai pour toute sous-catégorie de $(\mathrm{Ann})$.)
\end{enumerate}

\page{40}
(We can take for $C$ any subcategory of $E$ which is \emph{generating} $E$, in
the sense earlier explained, and take $E \xrightarrow{f_{*}} \widehat{C}$ by
\struck{associating to each} $f_{*}(X) = (A \mapsto \operatorname{Hom}(A, X))$ ;
to say that this functor is \emph{conservative} means that $C$ is generating, and
this \struck{does imply} is equivalent with $f_{*}$ fully faithful, and implies
the existence of the left adjoint $f^{*}$ \ldots)

\textbf{Generators.} Useful in various contexts: existence of adjoints
(Cartesian, or otherwise) \struck{\ill{}} ; together with generators, existence
of cogenerators, and technically the fact that \struck{functors} contravariant
functors to Sets $(S \to)$ that transform $\varinjlim$ ($\varprojlim$) into
$\varprojlim$ are representable.

NB Recall that a topos has cohomotopy and homotopy invariants.

\textbf{Examples}
\begin{itemize}
\item $\mathrm{Top}(X)$ : $\mathrm{Top}(\emptyset)$ empty topos ;
$\mathrm{Top}(1)$ point (or final) topos.
\item $\mathrm{Top}(X, G)$ : $\mathrm{Top}(X, e) = \mathrm{Top}(X)$ ;
$\mathrm{Top}(*, G) = B_{G}$. $H^{i}(B_{G}, \ ) = H^{i}(G, \ )$

Here fiber functor of $B_{G}$ is forgetful functor. It has automorphisms
\ill{}, and it \ill{} generalizes $B_{G}$.
\item $\widehat{C}$ (from algebra)
\item $C^{\circ} \xrightarrow{\text{fully faithful}} \mathrm{Fib}(\widehat{C})$
\quad [we can prove this is even $\mathrm{Ind}(C^{\circ}) \simeq
\mathrm{Fib}(\widehat{C})$]
\end{itemize}

\textbf{Question.} Can we have an equivalence
$\widehat{C} \simeq \mathrm{Top}(X)$ \struck{\ill{}} ? It would suffice that $C$
is an ordered category, \struck{\ill{}} and that $X$ has \ill{} which have a
smallest neighbourhood (« essential points ») --- they would be
\struck{everywhere} dense --- $C$ could be interpreted as contained in the set of
\ill{} of $X$ with the \ill{} \ill{} relation \ill{} we can then take
$X = \mathrm{Ind}(C)$ (dual to $C$ \ldots\ of decreasing families) and a
\ill{} the topology \ill{}

\marginal{\ill{} \ill{} it is \ill{}, and their \ill{} that
\ill{} $\Rightarrow$ different \ill{} $\Rightarrow$ fully faithful. We will
\ill{} that \ill{} categories \ill{} equivalent \ill{} topos \ill{}}

\end{document}
